\section{CONCLUSIONS}

\subsection{Summary of Findings}
This project set out to build, evaluate and visualise an end-to-end reinforcement-learning system for optimising web-advertisement selection, and to bridge the gap between the formal theory of multi-armed bandit algorithms and their observable, debuggable behaviour on a realistic advertising-style reward structure. We constructed an interactive, browser-based simulation framework that implements five canonical decision policies --- a Random baseline, $\varepsilon$-Greedy, Decaying $\varepsilon$-Greedy, the Upper Confidence Bound algorithm (UCB1), and Thompson Sampling --- against a unified agent interface, and runs them synchronously on a shared environment whose reward function can be configured as either binary click-through-rate or per-click revenue. The five algorithms are implemented from first principles in TypeScript with no third-party reinforcement-learning dependency, which makes the system fully auditable and dependency-light.

The experimental campaign described in Section~4 yielded three principal findings. First, principled bandit algorithms substantially outperform the Random baseline on the default scenario: cumulative regret over $T = 500$ steps is reduced from $97.4$ for Random to $8.7$ for UCB1, an order-of-magnitude improvement, accompanied by a near-quadrupling of the share of impressions allocated to the optimal arm (from $25\%$ to $96\%$). Second, UCB1 is the strongest performer on this scenario across every metric --- cumulative reward, regret, convergence step, and optimal-arm share --- and is robust to the choice of reward signal because its sample-mean update rule operates directly on the observed reward. Third, Thompson Sampling, despite its strong empirical reputation in the Bernoulli setting, under-performs $\varepsilon$-Greedy on the default Revenue-mode scenario, not because of an implementation defect but because its Beta--Bernoulli model is mis-specified for the revenue reward. This third finding, in our view, is the most consequential operational lesson of the project: the choice of algorithm and the choice of reward signal are joint decisions, not independent ones.

The sensitivity analyses (Section~4.3.8) confirmed that the ranking is robust. Across 100 randomly generated scenarios, UCB1 was the lowest-regret agent in $71\%$ of runs; $\varepsilon$-Greedy in $17\%$; Thompson Sampling in $10\%$ (predominantly on scenarios in which the CTR-optimal and revenue-optimal arms happened to coincide); and Decaying $\varepsilon$-Greedy and Random in $1\%$ each. UCB1's median regret across all 100 scenarios was a factor of 3.4 lower than $\varepsilon$-Greedy's, and the relative ranking was preserved when the arm count was varied from $K=2$ to $K=8$ and when mild non-stationarity was introduced into the environment.

Beyond the experimental results, the project delivers an artefact --- the simulator itself --- that operates simultaneously as a research benchmark (algorithmically rigorous, headlessly reproducible), as a teaching tool (every agent's internal state is visualised in real time), and as a reference implementation pattern (clean separation of environment, agent, simulation engine, store and presentation, with no dependency on a third-party RL library). The History view's per-algorithm win-rate analytics, computed across the 100 randomly generated scenarios, corroborate the headline ranking and provide a robust foundation against which future algorithms can be benchmarked. These outcomes provide a strong foundation for developing trustworthy, reinforcement-learning-assisted decision tools capable of supporting more informed ad-serving decisions and improved revenue generation for online advertising platforms.

\subsection{Limitations of the Work}
Despite the promising results, several limitations inherent to the scope and design of this project must be acknowledged. Each of these limitations represents a deliberate scope decision rather than an oversight; they are recorded here so that future work can address them systematically.

\paragraph{Synthetic environments only.} All experiments were run on synthetic Bernoulli environments. While this is the standard methodology in the bandit literature and is essential for ground-truth regret computation, it means the present results do not automatically transfer to real ad-serving traffic. Real CTRs are typically below one per cent, exhibit strong non-stationarity, and depend heavily on contextual features (user, page, time of day, device) that our context-free MAB formulation does not consume. The simulator can be configured to operate on low-CTR scenarios and to apply sinusoidal drift, partially mitigating these limitations, but a full validation on real logged data would require off-policy evaluation methodology that is outside the present scope.

\paragraph{No contextual features.} The current implementation is a non-contextual bandit: every arm has a fixed expected reward distribution independent of the user or the page. Real ad systems are almost universally contextual, and contextual bandits (LinUCB, Bayesian contextual variants, contextual-Thompson) offer substantially better performance when relevant features are available. Adding contextual features would require non-trivial changes to the agent interface and to the visualisation layer (per-context value estimates rather than per-arm). This is the single most impactful direction for future enhancement.

\paragraph{Limited treatment of non-stationarity.} The environment supports a sinusoidal drift mode, but none of the five implemented agents are designed specifically for non-stationary rewards. Sliding-window UCB, discounted UCB, and Thompson Sampling with a forgetting factor are well-established mitigations that would close most of the regret gap our experiments observed under drift. They were not implemented in the current scope for time reasons.

\paragraph{No real ad-serving integration.} The simulator does not connect to any real ad-server, real-time-bidding system, or production data source. It is therefore a research and teaching artefact, not a production tool. Integrating with a real ad infrastructure would require addressing safety (an explore action on a live user costs real money), latency (selection decisions in production are made in sub-10\,ms), and governance (which features fed which decision, for audit and compliance).

\paragraph{Limited statistical comparison framework.} The aggregate-comparison metrics reported in Section~4 are mean and standard deviation across 100 runs. A more rigorous comparison would include paired-difference tests (e.g.\ a Wilcoxon signed-rank test on the per-run regret differences between UCB1 and Thompson Sampling) and bootstrap confidence intervals on the win rates. These analyses are straightforward to add in a future iteration but were not included in the present submission.

\paragraph{Single horizon length tuned.} The hyperparameters for each algorithm (the decay rate $\beta$ for Decaying $\varepsilon$-Greedy, the confidence constant $c$ for UCB1, the prior $(\alpha_0, \beta_0)$ for Thompson Sampling) were not formally tuned for every horizon length we report. As discussed in Section~4.3.3, the decay schedule for Decaying $\varepsilon$-Greedy is mismatched to the $T = 500$ horizon; on longer horizons its ranking among the agents improves.

\paragraph{No formal external validation.} The benchmark is internal to the team. No external research group has independently re-run the experiments or audited the implementation. Publishing the code and inviting peer replication would strengthen the credibility of the empirical findings.

\paragraph{User-interface accessibility.} While the Radix UI primitives used in the implementation provide a strong baseline of ARIA support, a formal accessibility audit has not yet been performed. Internationalisation, right-to-left language support and screen-reader testing are open issues.

\subsection{Scope for Future Enhancements}
The findings of this project and the limitations above suggest several concrete directions for future enhancement, ordered roughly by anticipated impact.

\paragraph{Contextual bandits.} The natural next step is to extend the agent interface from \texttt{selectArm()} to \texttt{selectArm(context)} and to implement LinUCB and contextual Thompson Sampling with a linear posterior. The environment would correspondingly need to accept context vectors and produce per-context arm rewards. This change alone would more than double the practical relevance of the simulator and would bring it within striking distance of the algorithms that drive real production ad systems.

\paragraph{Sliding-window and forgetting variants.} Adding sliding-window UCB and discounted UCB (parameter $\gamma$) to the agent pool, together with a Thompson Sampling variant that decays the Beta parameters at each step, would let users compare stationary and non-stationary algorithms head-to-head on the drift environment. This is a high-value, low-effort extension and would close most of the regret gap observed under non-stationarity in Section~4.3.8.

\paragraph{Deep reinforcement-learning agents.} A Deep Q-Network agent is partially scaffolded in the codebase (\texttt{src/rl/agents/DQN.ts}) but is not currently wired into the agent-selection UI or the simulation engine. Completing the integration --- and validating that DQN does not under-perform UCB1 on the non-contextual default scenario, where it has no informational advantage --- would extend the simulator's pedagogical scope into the deep-RL family and provide a natural bridge from bandits to full reinforcement learning.

\paragraph{Off-policy evaluation toolkit.} For users who want to evaluate the implemented algorithms on real logged data, an off-policy evaluation module --- inverse-propensity-score (IPS) estimators, doubly-robust estimators, and self-normalised importance sampling --- would be a natural extension. The output would be an estimate, with confidence intervals, of what the cumulative reward of each implemented agent \emph{would have been} had it been deployed on the logged trajectory.

\paragraph{Multi-objective optimisation.} Real ad platforms balance multiple objectives: revenue, CTR, advertiser satisfaction, user retention. Extending the simulator to multi-objective bandits with explicit Pareto fronts, or constrained bandits with budget and fairness constraints, would expose the trade-offs more honestly than the current single-reward formulation.

\paragraph{Rigorisation of the comparison framework.} Adding statistical hypothesis testing (paired Wilcoxon, bootstrap CIs), explicit power analyses on the win-rate observations, and confidence-bounded comparison plots would lift the experimental campaign from informal to publication-grade.

\paragraph{Accessibility and internationalisation.} A formal accessibility audit using a screen reader and a keyboard-only navigation walkthrough, plus internationalisation infrastructure (translation files, locale-aware number formatting, right-to-left layout), would broaden the audience for the educational components of the system.

\paragraph{Mobile experience.} The current layout is optimised for desktop screens. A responsive redesign that collapses the three-pane simulator layout into a tabbed mobile experience, while preserving the live-update properties, would extend reach to phone and tablet users and make the educational artefact accessible in classrooms that cannot rely on desktops.

\paragraph{Integration with a real ad-exchange via mock.} Even short of a full production integration, a stub bridge that imports a sanitised, anonymised sample of click logs and lets the user run the implemented agents on the replayed traffic would be a substantial step toward real-world validity. The simulator's existing seed-replay infrastructure makes this straightforward to add.

\paragraph{Algorithm-comparison portal.} The win-rate analytics in the History view are currently visible only to the local user. Publishing them, anonymised and aggregated, to a small shared leaderboard would let multiple researchers contribute results on the same scenarios and create a small but useful open benchmark for bandit algorithms.

\subsection{Broader Impact and Ethical Considerations}
Reinforcement-learning systems applied to advertisement selection sit at a sensitive intersection of commercial, technical and ethical concerns. The same algorithmic improvements that increase platform revenue can, depending on the reward signal chosen, also amplify behaviour that is harmful to users or to the broader information ecosystem. This subsection records the considerations we believe a practitioner deploying such a system should keep front-of-mind.

\paragraph{Choice of reward signal as a value statement.}
Section~4.3.7 made the point that the reward signal is not a neutral technical detail: it is the objective the policy will maximise, and it determines which user behaviours are reinforced. A CTR-mode policy will reward whatever creative drives clicks, including clickbait. A revenue-mode policy will reward whatever creative converts into spend, which may favour predatory advertisers. A long-term-value-mode policy~[6] will reward retention and satisfaction at the expense of immediate revenue. The decision among these three is not a tuning choice; it is a policy choice that should involve product, legal and ethics stakeholders, not just engineering.

\paragraph{Manipulation and dark patterns.}
A sufficiently capable bandit algorithm will learn to optimise whatever target it is given. When the target is short-term engagement, the algorithm may converge on creatives that exploit known psychological levers (urgency, social proof, fear of missing out) without any awareness that this is what it is doing. Mitigating this requires upstream guard-rails on the creative inventory (creative-review policies, banned-content lists), not algorithmic constraints downstream. Our simulator does not currently include such guard-rails, but the same architectural seam --- the \texttt{AdEnvironment} class --- is the right place to add them.

\paragraph{Fairness across advertisers.}
A bandit that converges aggressively (low regret) by definition concentrates impressions on a small number of winning creatives. Concentrated impression share has competitive consequences for advertisers: a new entrant's creative may be denied the impressions it needs to demonstrate its quality because the bandit has already locked onto a known winner. UCB1's optimism-under-uncertainty mechanism partly mitigates this --- a new arm has an infinite confidence bonus on its first pull, guaranteeing at least one impression --- but in long-running systems with creative churn, additional fairness constraints (minimum impression share per advertiser, ramp-up periods for new creatives) are advisable. These are policy layers that sit \emph{on top of} the bandit, not replacements for it.

\paragraph{Privacy and data minimisation.}
The basic non-contextual MAB formulation studied in this project does not consume user-identifying data. Each step's information is a single arm pull and a single binary click outcome; no user features cross the decision boundary. This is one of the structural privacy advantages of context-free bandits compared to contextual bandits or deep RL-based personalisation, which typically consume browsing history, demographic features and behavioural signals. Should the system be extended to contextual bandits (Section~5.3), data-minimisation principles and the legal frameworks (GDPR, India's DPDP Act, CCPA) become directly relevant. Even with our context-free implementation, the per-run histories stored in browser local storage are user-visible and user-deletable from within the application UI; no logging or telemetry is transmitted off-device.

\paragraph{Transparency and auditability.}
A core design goal of the simulator is that every decision an agent makes is decomposable into the algorithm's current internal state plus the algorithm's selection rule. UCB1 selects the arm with the highest $Q(a) + c\sqrt{\ln t / N(a)}$; the user can see $Q$, $c$, $\ln t$ and $N$ at every step. Thompson Sampling samples from a Beta posterior; the alpha and beta of that posterior are surfaced in the UI. This level of transparency is exceeded by virtually no production ad system and is not strictly necessary for performance, but we believe it is an important property of educational and audit-grade systems. Production deployments aiming for similar transparency can use the same architectural pattern --- exposing internal quantities to a privileged debug interface --- without needing to expose them to end users.

\paragraph{Environmental sustainability.}
The simulator runs entirely client-side; no GPU or server inference is required. A typical 500-step run consumes a small, bounded amount of CPU on the user's device and no network traffic after the initial page load. By contrast, deep-RL bandit variants typically require GPU servers for training and inference, with associated energy and carbon costs~[31]. Within the family of bandit algorithms our simulator implements, all five are $\mathcal{O}(K)$ per step in both compute and memory, making them effectively free in environmental terms. This is a real advantage of classical multi-armed bandits over their deep-RL cousins, and one that deserves more weight when production systems are designed.

\subsection{Lessons Learned}
Beyond the technical findings, the project produced four practical lessons that the team will carry into future work.

\paragraph{Clean interfaces enable parallel work.}
Once the \texttt{BaseAgent} contract was specified --- two methods, \texttt{selectArm} and \texttt{update}, plus shared statistics --- the five agent implementations, the environment, the simulation engine and the chart components could all be developed in parallel by different team members and integrated with almost no re-work. The lesson is generalisable: spending the first day specifying the interface, before any production code is written, pays for itself many times over in a multi-person project.

\paragraph{Persisting seeds is the easiest form of reproducibility.}
The simulator's history store records the random seed of every saved run. This single design choice turned what would otherwise have been "trust the screenshot" into "re-run the experiment yourself in 200 ms." Several times during the project, a reviewer's question about a specific number was answered by re-running the exact scenario from the History view. Reproducibility of this form costs almost nothing to implement and we will adopt it in all future experimental work.

\paragraph{Live visualisation is a debugging tool, not a cosmetic feature.}
Several subtle bugs in the early agent implementations were spotted only when their internal state was rendered live in the dashboard. The most memorable was a sign error in the UCB1 confidence bonus that produced the wrong selection on a small fraction of steps; the aggregate metrics still looked plausible, but the live UCB-score visualisation showed the bonus shrinking in the wrong direction relative to the per-arm counts. We have since adopted the principle that any non-trivial algorithmic state should be inspectable in real time during development.

\paragraph{Beware proxy objectives.}
The Revenue-mode mismatch finding (Section~4.3.7) was the single most surprising result of the project, and it was discovered almost by accident: we built the Revenue mode mainly to have a more interesting reward distribution for the visualisation, and only afterward noticed that Thompson Sampling was failing to find the right arm. The general lesson is that proxy objectives (CTR-as-stand-in-for-revenue) can systematically mislead even mathematically well-behaved algorithms, and that simulation environments which expose the gap between proxy and true objective have substantial pedagogical value.

In closing, this project demonstrates that the multi-armed bandit family --- a body of theory dating back to Robbins (1952) and Thompson (1933) --- remains directly relevant to one of the largest practical decision problems in modern computing. By implementing the canonical algorithms from first principles, instrumenting them with rich introspection, and exposing them through a transparent and pedagogically accessible interface, we hope the system serves both as a benchmark for future bandit research and as a teaching tool for the next generation of practitioners who will design the advertising and recommendation systems of the open web.
\newpage
